\documentclass[aps,preprint]{revtex4}%
\usepackage{amsfonts}
\usepackage{amsmath}
\usepackage{amssymb}
\usepackage{graphicx}%
\setcounter{MaxMatrixCols}{30}
%TCIDATA{OutputFilter=latex2.dll}
%TCIDATA{Version=5.50.0.2953}
%TCIDATA{CSTFile=revtex4.cst}
%TCIDATA{Created=Thursday, March 20, 2008 15:20:09}
%TCIDATA{LastRevised=Wednesday, April 09, 2008 12:02:23}
%TCIDATA{<META NAME="GraphicsSave" CONTENT="32">}
%TCIDATA{<META NAME="SaveForMode" CONTENT="1">}
%TCIDATA{BibliographyScheme=BibTeX}
%TCIDATA{<META NAME="DocumentShell" CONTENT="Articles\SW\REVTeX 4">}
%TCIDATA{Language=American English}
%BeginMSIPreambleData
\providecommand{\U}[1]{\protect\rule{.1in}{.1in}}
%EndMSIPreambleData
\newtheorem{theorem}{Theorem}
\newtheorem{acknowledgement}[theorem]{Acknowledgement}
\newtheorem{algorithm}[theorem]{Algorithm}
\newtheorem{axiom}[theorem]{Axiom}
\newtheorem{claim}[theorem]{Claim}
\newtheorem{conclusion}[theorem]{Conclusion}
\newtheorem{condition}[theorem]{Condition}
\newtheorem{conjecture}[theorem]{Conjecture}
\newtheorem{corollary}[theorem]{Corollary}
\newtheorem{criterion}[theorem]{Criterion}
\newtheorem{definition}[theorem]{Definition}
\newtheorem{example}[theorem]{Example}
\newtheorem{exercise}[theorem]{Exercise}
\newtheorem{lemma}[theorem]{Lemma}
\newtheorem{notation}[theorem]{Notation}
\newtheorem{problem}[theorem]{Problem}
\newtheorem{proposition}[theorem]{Proposition}
\newtheorem{remark}[theorem]{Remark}
\newtheorem{solution}[theorem]{Solution}
\newtheorem{summary}[theorem]{Summary}
\newenvironment{proof}[1][Proof]{\noindent\textbf{#1.} }{\ \rule{0.5em}{0.5em}}
\begin{document}
\preprint{ }
\title{Some Simple Explicit Examples of Correlated One Electron Theory and Density
Matrix Functionals }
\author{B. Weiner}
\affiliation{Department of Physics, Pennsylvania State University, DuBois PA 15801}

\pacs{}

\begin{abstract}
Some examples are presented displaying how General Spin Orbitals (GSO's) and
their occupation numbers uniquely determine N-electron Antisymmetrized Geminal
Power states. This provides a correlated one electron model that generalizes
the independent particle model that is universally used in picturing chemical
behavior. These examples also demonstrate that one electron models, involving
GSO's and occupation numbers, and ones based on a first order density matrix
are not equivalent. This difference is shown by displaying sets of different
GSO's that have the same occupation numbers that correspond to the same
density matrix but different N-electron states. This observation also leads to
the explicit construction of a density matrix functional.

\end{abstract}
\date{April 9th, 2008}
\startpage{1}
\maketitle


\section{Introduction\label{Intro}}

The utility and power of pictorial models in the understanding and
manipulation of complex physical phenomena is clearly evident in all areas of
human scientific endeavour. Perhaps one of the most ubiquitous examples of
this is the use of molecular orbitals in understanding structure, properties
and reactivity in chemistry and material science. This efficacy is due to the
approximate validity of the Independent Particle State (IPS) model, where

\begin{itemize}
\item the properties of individual electrons and their effective interactions
\emph{determine uniquely} the many electron state and its properties

\item electrons can be identified with orbitals and their observable
quantities (e.g. energy, ionization potential, electron affinity etc..) in a
one to one fashion

\item orbitals provide a geometric representation of electrons and thus multi
electron states.
\end{itemize}

The IPS model, however, can predict incorrect properties, reactions and
structures due to the absence of electron-electron correlation effects. In
order to overcome this defect more accurate methods such as Configuration
Interaction, Coupled Cluster, Perturbation Theory and many others have been
utilized that contain an explicit description of these interactions. Although
the resulting improved multi electron states still determine sets of orbitals,
now greater than the number of electrons and weighted with occupation numbers,
through the First Order Reduced Density Operator (FORDO), these orbitals and
their occupation numbers \emph{do not}, in general, determine a unique multi
electron state. Thus one looses the concrete pictorial aspects of a one
electron theory, as the multi electron state \emph{cannot} be described solely
in terms of single electron quantities and geometric forms.

Maintaining a one electron picture while including electron correlation
effects, i.e. constructing a correlated one electron picture, is a central
driving force in Kohn-Sham Density Functional Theory (KSDFT), Density Matrix
Functional Theory (DMFT) and various quasi particle approaches. A common
drawback in all of these methods is that it is not known how to express the
multi electron state as an explicit functional of orbitals and occupation
numbers. However, there is a class of correlated multi electron states and
hence their energies that can be expressed exactly as known functionals of
orbitals and occupation numbers, they are the Antisymmetrized Geminal Power
(AGP) States \cite{OrtizWeiner2006}.

The analysis of AGP states also clearly shows that a model based on orbitals
and occupation numbers is not equivalent to one based on FORDO's, as many sets
of orbitals determine the same FORDO while determining different $N$-electron
states. This observation allows one to explicitly construct examples of DMF's
whose domain is restricted to AGP's.

In this paper we review the definition of a correlated one particle theory
(section 2), the structure and parameterization of AGP\ states, their
$N$-electron energy functional, the expression of their energy in terms of one
electron quantities and discuss how one can interpret them pictorially
(section 3). In section 4 we define the AGP\ DMF and in section 5 we describe
simple examples of (i) AGP states that are determined by orbitals and
occupation numbers and (ii) AGP DMF's.

\section{One Electron Theory\label{OET}}

In order to be unambiguous we introduce the following definition: a One
Electron Theory (OET) is one in which the

\begin{itemize}
\item $N$-electron state is \emph{totally determined} \ by a set of
orthonormal Generalized Spin Orbitals (GSO's), (see Appendix \ref{Spin}),
$\left\{  \left\vert \psi_{j}\right\rangle ;1\leq j\leq r\right\}  $ spanning
one particle space and a set of occupation numbers $\left\{  N_{j};1\leq j\leq
r\right\}  $ of these GSO's

and

\item $N$-electron state energy can be expressed in terms of one electron
observables e.g. GSO kinetic energies, ionization potentials etc...
\end{itemize}

A OET is \emph{distinct} from ones based on FORDO's, as different sets of
GSO's (with the same occupation numbers) can produce the \emph{same} FORDO
while forming different $N$-electron states, with different energies. This
difference is due to the fact that the FORDO is invariant to the action of the
unitary group $U(1)$ on one electron state space, while in general the
$N$-electron state density operator is not. The $U(1)$ invariance can be
observed by considering the natural expansion of the FORDO
\begin{equation}
D^{1}=\sum_{1\leq j\leq r}N_{j}\left\vert \psi_{j}\right\rangle \left\langle
\psi_{j}\right\vert ,
\end{equation}
which is invariant to the action
\begin{align}
U(1)\left(  \left\vert \psi_{1}\right\rangle ,\cdots,\left\vert \psi
_{r}\right\rangle \right)   &  =\left(  \left\vert \psi_{1}\right\rangle
,\cdots,\left\vert \psi_{r}\right\rangle \right)  \left(
\begin{array}
[c]{ccc}%
e^{i\theta_{1}} & \cdots & 0\\
\vdots & \ddots & \vdots\\
0 & \cdots & e^{i\theta_{r}}%
\end{array}
\right) \nonumber\\
&  =\left(  e^{i\theta_{1}}\left\vert \psi_{1}\right\rangle ,\cdots
,e^{i\theta_{r}}\left\vert \psi_{r}\right\rangle \right)
\end{align}
of the unitary group $U(1)$ on the GSO's as%
\begin{equation}
D^{1}=\sum_{1\leq j\leq r}N_{j}e^{i\theta_{j}}\left\vert \psi_{j}\right\rangle
\left\langle \psi_{j}\right\vert e^{-i\theta_{j}}=\sum_{1\leq j\leq r}%
N_{j}\left\vert \psi_{j}\right\rangle \left\langle \psi_{j}\right\vert .
\end{equation}
The non invariance of the $N$-electron state density operator will be shown by
simple examples in the proceeding, where we show explicitly that Second Order
Reduced Density Operator's (SORDO's) of AGP\ states (thus any higher order
density operators) are not in general invariant to such transformations.

The observation that many sets of GSO's map into the same FORDO while
producing different $N$-electron states allows one to construct concrete
examples of DMF's (which are functionals of matrix representations of FORDO's)
through a Levy-Lieb \cite{Levy1979, Lieb1983}\ type of construction. Simple
examples of such functionals will also be displayed in the proceeding.

\section{Antisymmetrized Geminal Power States\label{AGP States}}

The most general form of an AGP state is produced by forming an
antisymmetrized power of a single geminal composed of GSO's. Geminals can be
expanded uniquely in parametric form as
\begin{equation}
\left\vert g\left(  \boldsymbol{c},\boldsymbol{\psi},\boldsymbol{\theta
}\right)  \right\rangle =\sum_{1\leq j\leq s}c_{j}e^{i\theta_{j}}\left\vert
\psi_{j}\psi_{j+s}\right\rangle ,\label{gem}%
\end{equation}
where $\left\{  \left\vert \psi_{j}\right\rangle ;1\leq j\leq r=2s\right\}  $
is a complete orthonormal basis (conb) of GSO's, the canonical coefficients
$\mathbf{c}$ lie on the positive portion of the hypersphere
\begin{equation}
\sum\limits_{1\leq j\leq s}c_{j}^{2}=1;\,c_{j}\geq0;1\leq j\leq s,
\end{equation}
the angles $\boldsymbol{\theta}=\left(  \theta_{1},\cdots,\theta_{s}\right)  $ satisfy%

\begin{equation}
0\leq\theta_{2},\cdots,\theta_{s}<2\pi;\,\theta_{1}=0 \label{Phases}%
\end{equation}
and the GSO's $\left\{  \psi_{j},\psi_{j+s}\right\}  $ belong to a set of non
equivalent pairs, $\mathcal{P},$ (see Appendix \ref{GSOpairs} and subsection
\ref{DenOp}).

\subsection{Density Operators\label{DenOp}}

The FORDO\ of an AGP\ is always at least doubly degenerate i.e.
\begin{equation}
D^{1}\left(  \boldsymbol{c},\boldsymbol{\psi}\right)  =\sum_{1\leq j\leq
s}N_{j}\left(  \boldsymbol{c}\right)  \left\{  \left\vert \psi_{j}%
\right\rangle \left\langle \psi_{j}\right\vert +\left\vert \psi_{j+s}%
\right\rangle \left\langle \psi_{j+s}\right\vert \right\}  ,\label{D1}%
\end{equation}
and it is uniquely parameterized by $\left\{  \boldsymbol{c},\boldsymbol{\psi
}\right\}  ,$ where $\boldsymbol{\psi}\mathbf{\in}\mathcal{P}$ are
representatives of the equivalence classes, $\left[  \boldsymbol{\psi}\right]
$, of paired normalized GSO's, which for non degenerate AGP's are defined by
\begin{equation}
\left\{  \psi_{a},\psi_{b}\right\}  \sim\left\{  \psi_{c},\psi_{d}\right\}
\Longleftrightarrow\left\vert \psi_{a}\psi_{b}\right\rangle \left\langle
\psi_{a}\psi_{b}\right\vert =\left\vert \psi_{c}\psi_{d}\right\rangle
\left\langle \psi_{c}\psi_{d}\right\vert ,\label{Equival}%
\end{equation}
for degenerate AGP's the equivalence is more involved and is discussed in
Appendix \ref{GSOpairs}. The occupation numbers $\left\{  N_{j}\left(
\boldsymbol{c}\right)  ;1\leq j\leq s\right\}  $ are given in terms of
$\left\{  n_{j};1\leq j\leq s\right\}  $ in a one-one fashion \cite{Yukalov00}
by
\begin{equation}
N_{j}\left(  \boldsymbol{c}\right)  =n_{j}S_{\frac{N}{2}-1}\left(  \hat
{\jmath}\right)  \left(  S_{\frac{N}{2}}\right)  ^{-1},\label{OccNum}%
\end{equation}
where $n_{j}=c_{j}^{2}$
\begin{equation}
S_{_{\frac{N}{2}}}\equiv S_{_{\frac{N}{2}}}\left(  \boldsymbol{n}\right)
=\sum_{1\leq j_{1}<\cdots<j_{\frac{N}{2}}\leq s}n_{j_{1}}\cdots n_{j_{_{\frac
{N}{2}}}},
\end{equation}
and
\begin{equation}
S_{\frac{N}{2}-1}\left(  \hat{\jmath}\right)  =\sum_{\substack{1\leq
j_{1}<\cdots<j_{\frac{N}{2}-1}\leq s\\j_{1},\cdots,j_{\frac{N}{2}-1}\neq
j}}n_{j_{1}}\cdots n_{j_{_{\frac{N}{2}-1}}}.
\end{equation}
The SORDO\ is given by
\begin{align}
&  D^{2}\left(  \boldsymbol{c},\boldsymbol{\psi},\boldsymbol{\theta}\right)
=\left(  S_{\frac{N}{2}}\right)  ^{-1}\left\{  \sum_{1\leq j\leq s}%
n_{j}S_{\frac{N}{2}-1}\left(  \hat{\jmath}\right)  \left\vert \psi_{j}%
\psi_{j+s}\right\rangle \left\langle \psi_{j}\psi_{j+s}\right\vert +\right.
\nonumber\\
&  \sum_{1\leq j\neq k\leq s}c_{j}c_{k}e^{i\left(  \theta_{j}-\theta
_{k}\right)  }S_{\frac{N}{2}-1}\left(  \hat{\jmath}\hat{k}\right)  \left\vert
\psi_{j}\psi_{j+s}\right\rangle \left\langle \psi_{k}\psi_{k+s}\right\vert
\left.  +\sum_{\substack{1\leq j<k\leq r\\k\neq j+s}}n_{j}n_{k}S_{\frac{N}%
{2}-2}\left(  \hat{\jmath}\hat{k}\right)  \left\vert \psi_{j}\psi
_{k}\right\rangle \left\langle \psi_{j}\psi_{k}\right\vert \right\}  ,
\end{align}
where
\begin{align}
S_{\frac{N}{2}-1}\left(  \hat{\jmath}\hat{k}\right)   &  =\sum
_{\substack{1\leq j_{1}<\cdots<j_{\frac{N}{2}-1}\leq s\\j_{1},\cdots
,j_{\frac{N}{2}-1}\neq j,k}}n_{j_{1}}\cdots n_{j_{_{\frac{N}{2}-1}}},\\
S_{\frac{N}{2}-2}\left(  \hat{\jmath}\hat{k}\right)   &  =\sum
_{\substack{1\leq j_{1}<\cdots<j_{\frac{N}{2}-2}\leq s\\j_{1},\cdots
,j_{\frac{N}{2}-1}\neq j,k}}n_{j_{1}}\cdots n_{j_{_{\frac{N}{2}-2}}},\\
S_{0}\left(  \hat{\jmath}_{1}\cdots\hat{\jmath}_{p}\right)   &  =1;\quad
S_{m}\left(  \hat{\jmath}_{1}\cdots\hat{\jmath}_{p}\right)  =0,\,\quad
m<0;\quad0\leq p\leq s.
\end{align}


\subsection{State Vector\label{StateVec}}

The $N$-electron state vector has the CI form
\begin{equation}
\left\vert g\left(  \boldsymbol{c},\boldsymbol{\psi},\boldsymbol{\theta
}\right)  ^{\frac{N}{2}}\right\rangle =\sum_{1\leq j_{1}<\cdots<j_{\frac{N}%
{2}}\leq s}c_{j_{1}}\cdots c_{j_{\frac{N}{2}}}\exp\left\{  i\sum
\limits_{1\leq\kappa\leq\frac{N}{2}}\theta_{j_{\kappa}}\right\}  \left\vert
\psi_{j_{1}}\psi_{j_{1}+s}\cdots\psi_{j_{\frac{N}{2}}}\psi_{_{j_{\frac{N}{2}}%
}+s}\right\rangle ,
\end{equation}
which explicitly displays the dependence of AGP states, that have the same
FORDO (those with common $\left(  \boldsymbol{c},\boldsymbol{\psi}\right)  $),
on the phases $\left\{  \left.  \theta_{j_{\kappa}}\right\vert 1\leq\kappa\leq
s\right\}  $.

\subsection{Energy Functional\label{AGPFunc}}

The energy of an AGP\ state is given by
\begin{align}
&  \mathcal{E}\left(  \boldsymbol{c},\boldsymbol{\psi},\boldsymbol{\theta
}\right)  =\frac{\mathcal{H}\left(  \boldsymbol{c},\boldsymbol{\psi
},\boldsymbol{\theta}\right)  }{S_{\frac{N}{2}}\left(  \boldsymbol{n}\right)
}=\nonumber\\
&  \left\{  \sum_{1\leq j\leq s}n_{j}S_{\frac{N}{2}-1}\left(  \hat{\jmath
}\right)  \left\{  \overset{}{\left\langle \psi_{j}\left\vert h_{1}\psi
_{j}\right.  \right\rangle +\left\langle \psi_{j+s}\left\vert h_{1}\psi
_{j+s}\right.  \right\rangle }+\,\left\langle \psi_{j}\psi_{j+s}\right\vert
\left\vert \psi_{j}\psi_{j+s}\right\rangle \right\}  \right.  \nonumber\\
&  \qquad\qquad\qquad+\sum_{1\leq j\neq k\leq s}\left\{  c_{j}c_{k}e^{i\left(
\theta_{j}-\theta_{k}\right)  }S_{\frac{N}{2}-1}\left(  \hat{\jmath}\hat
{k}\right)  \left\langle \psi_{j}\psi_{j+s}\right\vert \left\vert \psi_{k}%
\psi_{k+s}\right\rangle \right.  \nonumber\\
&  \qquad\qquad\qquad\qquad\left.  \left.  \left.  +\,\frac{1}{2}\overset
{}{n_{j}n_{k}S_{\frac{N}{2}-2}\left(  \hat{\jmath}\hat{k}\right)  \left\langle
\psi_{j}\psi_{k}|||\psi_{j}\psi_{k}\right\rangle }\right\}  \right\}  \right/
S_{\frac{N}{2}}\left(  \boldsymbol{n}\right)  .\label{EnFunc}%
\end{align}
where the terms involving the coulomb interaction are defined as
\begin{align}
\left\langle \psi_{j}\psi_{k}|||\psi_{j}\psi_{k}\right\rangle  &
=\left\langle \psi_{j}\psi_{k}||\psi_{j}\psi_{k}\right\rangle +\left\langle
\psi_{j}\psi_{k+s}||\psi_{j}\psi_{k+s}\right\rangle \nonumber\\
&  +\left\langle \psi_{j+s}\psi_{k}||\psi_{j+s}\psi_{k}\right\rangle
+\left\langle \psi_{j+s}\psi_{k+s}||\psi_{j+s}\psi_{k+s}\right\rangle
\end{align}
and one can express the one and two electron expectation values in terms of of
space-spin wave functions with coordinates $\left(  \boldsymbol{r}%
,\boldsymbol{\xi}\right)  $ in atomic units as
\begin{equation}
\left\langle \psi_{j}\left\vert h_{1}\psi_{j}\right.  \right\rangle =\int
\psi_{j}^{\ast}\left(  \boldsymbol{r},\boldsymbol{\xi}\right)  \left\{
-\frac{1}{2}\nabla_{\boldsymbol{r}}^{2}+V_{\text{ext}}\left(  \boldsymbol{r}%
\right)  \right\}  \psi_{j}\left(  \boldsymbol{r},\boldsymbol{\xi}\right)
d\boldsymbol{r}d\boldsymbol{\xi}%
\end{equation}
and
\begin{align}
&  \left\langle \psi_{j}\psi_{k}||\psi_{l}\psi_{m}\right\rangle =\int\left\{
\psi_{j}^{\ast}\left(  \boldsymbol{r}_{1},\boldsymbol{\xi}_{1}\right)
\psi_{k}^{\ast}\left(  \boldsymbol{r}_{2},\boldsymbol{\xi}_{2}\right)
\frac{1}{\left\Vert \boldsymbol{r}_{1}-\boldsymbol{r}_{2}\right\Vert }\psi
_{l}\left(  \boldsymbol{r}_{1},\boldsymbol{\xi}_{1}\right)  \psi_{m}\left(
\boldsymbol{r}_{2},\boldsymbol{\xi}_{2}\right)  \right.  \nonumber\\
&  -\left.  \psi_{j}^{\ast}\left(  \boldsymbol{r}_{1},\boldsymbol{\xi}%
_{1}\right)  \psi_{k}^{\ast}\left(  \boldsymbol{r}_{2},\boldsymbol{\xi}%
_{2}\right)  \frac{1}{\left\Vert \boldsymbol{r}_{1}-\boldsymbol{r}%
_{2}\right\Vert }\psi_{l}\left(  \boldsymbol{r}_{2},\boldsymbol{\xi}%
_{2}\right)  \psi_{m}\left(  \boldsymbol{r}_{1},\boldsymbol{\xi}_{1}\right)
\right\}  d\boldsymbol{r}_{1}d\boldsymbol{r}_{2}d\boldsymbol{\xi}%
_{1}d\boldsymbol{\xi}_{2}%
\end{align}
respectively.

\subsection{One Electron Picture\label{OneElecPic}}

One can canonically absorb the phase factors in eq.$\left(  \ref{gem}\right)
$ into the GSO's by setting
\begin{equation}
\left\vert g\left(  \boldsymbol{c},\boldsymbol{\psi},\boldsymbol{\theta
}\right)  \right\rangle =\sum_{1\leq j\leq s}c_{j}e^{i\theta_{j}}\left\vert
\psi_{j}\psi_{j+s}\right\rangle =\sum_{1\leq j\leq s}c_{j}\left\vert
\tilde{\psi}_{j}\tilde{\psi}_{j+s}\right\rangle ,
\end{equation}
where
\begin{align}
\left\vert \tilde{\psi}_{j}\right\rangle  &  =e^{i\theta_{j}/2}\left\vert
\psi_{j}\right\rangle \nonumber\\
\left\vert \tilde{\psi}_{j+s}\right\rangle  &  =e^{i\theta_{j}/2}\left\vert
\psi_{j+s}\right\rangle
\end{align}
showing that the geminal is \emph{uniquely defined} by the GSO's $\left\{
\left.  \left\vert \tilde{\psi}_{j}\right\rangle \right\vert 1\leq j\leq
r\right\}  $ and the occupation numbers $\left\{  \left.  N_{j}\right\vert
1\leq j\leq r\right\}  $ (through the inverse of the function defined by eq.
$\left(  \ref{OccNum}\right)  ),$ thus satisfying the first of the two
criteria for a OET displayed in section \ref{OET}.

\subsection{One Electron Quantities\label{OEP}}

In \cite{OrtizWeiner2006} it was shown that for AGP's states, $\left\vert
g\left(  \boldsymbol{c}_{\#},\boldsymbol{\tilde{\psi}}_{\#}\right)  ^{\frac
{N}{2}}\right\rangle $ that optimize the energy functional eq. $\left(
\ref{EnFunc}\right)  $ i.e.
\begin{equation}
\mathcal{E}\left(  \boldsymbol{c}_{\#},\boldsymbol{\tilde{\psi}}_{\#}\right)
=\min_{\boldsymbol{c},\boldsymbol{\tilde{\psi}}}\mathcal{E}\left(
\boldsymbol{c},\boldsymbol{\tilde{\psi}}\right)
\end{equation}
one can express the total $N$-electron energy in terms of the occupation
numbers $\left\{  \left.  N_{j}\left(  \boldsymbol{c}_{\#}\right)  \right\vert
1\leq j\leq r\right\}  $, Ionization Potentials (IP's), $\left\{  \left.
I\left(  \tilde{\psi}_{\#j}\right)  \right\vert 1\leq j\leq r\right\}  $ and
external potentials plus kinetic energies $\left\{  \left.  h_{1}\left(
\tilde{\psi}_{\#j}\right)  \right\vert 1\leq j\leq r\right\}  $ associated
with the optimized GSO's $\left\{  \left.  \left\vert \tilde{\psi}%
_{\#j}\right\rangle \right\vert 1\leq j\leq r\right\}  $. The tilde signifies
that the phase factor dependence has been absorbed into the GSO's as shown in
subsection \ref{OneElecPic}. The IP, $I\left(  \tilde{\psi}_{\#j}\right)  ,$
is the difference between the optimized $N$-electron state energy and the
$N$-$1$ electron state energy formed by the same GSO's as the $N$-electron one
but with the j$^{\text{th}}$ GSO\ removed
\begin{equation}
I\left(  \tilde{\psi}_{\#j}\right)  =\mathcal{E}\left(  \boldsymbol{c}%
_{\#},\boldsymbol{\tilde{\psi}}_{\#}\right)  -\mathcal{E}_{N\,-1}\left(
\boldsymbol{c}_{\#},\boldsymbol{\tilde{\psi}}_{\#},\tilde{\psi}_{\#j}\right)
\end{equation}
i.e. the frozen IP\ potential of the GSO $\left\vert \tilde{\psi}%
_{\#j}\right\rangle $. The one electron expectation value $h_{1}\left(
\tilde{\psi}_{\#j}\right)  $ is the sum of the kinetic and external
interaction energies and is given as
\begin{equation}
h_{1}\left(  \tilde{\psi}_{\#j}\right)  =\left\langle \left.  \tilde{\psi
}_{\#j}\right\vert h\tilde{\psi}_{\#j}\right\rangle =\left\langle \left.
\tilde{\psi}_{\#j}\right\vert \left(  K+V_{\text{ext}}\right)  \tilde{\psi
}_{\#j}\right\rangle .
\end{equation}
The expression for the total energy is
\begin{equation}
\mathcal{E}\left(  \boldsymbol{c}_{\#},\boldsymbol{\tilde{\psi}}_{\#}\right)
=\frac{1}{2}\sum_{1\leq j\leq2s}N_{j}\left(  \boldsymbol{c}_{\#}\right)
\kappa\left(  \tilde{\psi}_{\#j}\right)  ,\label{OETenergy}%
\end{equation}
where
\begin{equation}
\kappa\left(  \tilde{\psi}_{\#j}\right)  =I\left(  \tilde{\psi}_{\#j}\right)
+h_{1}\left(  \tilde{\psi}_{\#j}\right)  \label{kappa}%
\end{equation}
and $N_{j}\left(  \boldsymbol{c}_{\#}\right)  $ is the optimized occupation
number of the GSO\ $\left\vert \tilde{\psi}_{\#j}\right\rangle .$

\subsection{Orbital Models\label{OrbModel}}

The one electron independent particle model is based on GSO's obtained by
energy optimizing a single determinantal wave function, that have integer
occupation numbers. As the number of electrons and orthonormal GSO's $\left\{
\left.  \left\vert \psi_{j}\right\rangle \right\vert 1\leq j\leq N\right\}  $
that define the $N$-electron state in the IPS model are equal they can be
identified with the Probability Densities (PD's)
\begin{equation}
p_{\psi_{j}}\left(  \boldsymbol{r},\boldsymbol{\xi}\right)  =\psi_{j}^{\ast
}\left(  \boldsymbol{r},\boldsymbol{\xi}\right)  \psi_{j}\left(
\boldsymbol{r},\boldsymbol{\xi}\right)  \label{PD}%
\end{equation}
produced by these GSO's, which can then interpreted as charge clouds. These
GSO's are not unique as $N$-electron IPS's are invariant to the unitary group
$U\left(  N\right)  $ that map the $N$-dimensional subspace generated by
$\left\{  \left.  \left\vert \psi_{j}\right\rangle \right\vert 1\leq j\leq
N\right\}  $ to itself. Hence any set of orthonormal GSO's
\begin{equation}
\left\{  \left.  u\left\vert \psi_{j}\right\rangle \right\vert 1\leq j\leq
N\right\}  ;\quad u\in U\left(  N\right)  \label{UGSO}%
\end{equation}
and their associated PD's, \emph{which in general have different shapes}, give
equally valid descriptions of the distribution of the $N$-electrons. In order
to remove this pictorial ambiguity one can consider the optimized Fock
operator, which is represented in the basis of optimized GSO's by a matrix
which is divided into two blocks one $N\times N$ (the occupied GSO's) the
other $\left(  r-N\right)  \times\left(  r-N\right)  $ (the virtual GSO's).
Canonical occupied GSO's can then be defined as the eigen vectors of the
occupied block by finding a $u\in U\left(  N\right)  $ that diagonalizes this
block. The associated eigenvalues are ionization potentials associated with
these canonical GSO's. As both the FORDO of the optimized IPS and the state
itself are invariant to the group $U\left(  N\right)  $ the PD's produced by
the canonical GSO's can then be used to form a standard pictorial description
of the electronic distribution. However, from basic quantum mechanical
considerations the canonical GSO's and the resulting PD's are not any more
meaningful than any other set given by eq. $\left(  \ref{UGSO}\right)  $.
Other criteria, such as localization properties can and have been utilized
that produce different visualizations, which are physically equivalent, but
emphasize different aspects of the electronic structure. The importance of a
specific canonical GSO to the reactivity of a given material/molecule can be
estimated from the expression eq. $\left(  \ref{OETenergy}\right)  $, which
can be expanded as
\begin{equation}
\mathcal{E}\left(  \boldsymbol{\tilde{\psi}}_{\#}\right)  =\frac{1}{2}%
\sum_{1\leq j\leq N}\kappa\left(  \tilde{\psi}_{\#j}\right)  =\frac{1}{2}%
\sum_{1\leq j\leq N}\left\{  \epsilon\left(  \tilde{\psi}_{\#j}\right)
+h_{1}\left(  \tilde{\psi}_{\#j}\right)  \right\}  ,\label{FockEn}%
\end{equation}
where $\left\{  \left.  \epsilon\left(  \tilde{\psi}_{\#j}\right)  \right\vert
1\leq j\leq N\right\}  $ are the eigenvalues of the occupied GSO's of the Fock
operator. The highest energy GSO\ $\left\vert \tilde{\psi}_{\#N}\right\rangle
,$ where $\kappa\left(  \tilde{\psi}_{\#N}\right)  \geq\cdots\geq\kappa\left(
\tilde{\psi}_{\#1}\right)  ,$ corresponds to the most reactive "electron",
(note we have used the \emph{non conventional} association between orbitals
and energies given by $\left\{  \kappa\left(  \tilde{\psi}_{\#j}\right)
\right\}  $ which takes into account both the IP, effective coulombic
interaction and the kinetic energy and can be deduced by the use of the virial theorem).

The OET introduced in section \ref{OET} generalizes this picture to correlated
electrons described by AGP states, which contain IPS's as a special case. Now,
however, the $N$-electron state is determined by more than $N$ GSO's with non
integer occupation numbers and thus one cannot associate an individual
electron with a specific PD. It is still possible, however, to visualize the
structure in terms of electrons distributed according to the PD's produced by
the representative canonical orbitals \emph{of the optimized geminal}, but now
weighted with the associated occupation numbers. The total energy can again be
expressed as eq. $\left(  \ref{OETenergy}\right)  $ which can be recast in
terms of the diagonal elements $\left\{  \left.  \epsilon\left(  \tilde{\psi
}_{\#j}\right)  \right\vert 1\leq j\leq r\right\}  $ of a generalized Fock
operator\emph{ }\cite{OrtizWeiner2006} as
\begin{align}
\mathcal{E}\left(  \boldsymbol{c}_{\#,}\boldsymbol{\tilde{\psi}}_{\#}\right)
&  =\frac{1}{2}\sum_{1\leq j\leq s}N\left(  \tilde{\psi}_{\#j}\right)
\left\{  \kappa\left(  \tilde{\psi}_{\#j}\right)  +\kappa\left(  \tilde{\psi
}_{\#j+s}\right)  \right\}  \nonumber\\
&  =\frac{1}{2}\sum_{1\leq j\leq s}N\left(  \tilde{\psi}_{\#j}\right)
\left\{  \epsilon\left(  \tilde{\psi}_{\#j}\right)  +\epsilon\left(
\tilde{\psi}_{\#j+s}\right)  +h_{1}\left(  \tilde{\psi}_{\#j}\right)
+h_{1}\left(  \tilde{\psi}_{\#j+s}\right)  \right\}  .\label{GenFocken}%
\end{align}
The generalized Fock matrix is in general \emph{non diagonal} i.e.
$\epsilon\left(  \tilde{\psi}_{\#j}\right)  $ are not eigenvalues and is a
\emph{full} $r\times r$ matrix. The most reactive electron now corresponds to
PD's $\left\{  \left.  p_{j}\right\vert j\in\left\{  j_{1H},\cdots
,j_{mH}\right\}  \right\}  $, that are weighted by the occupation numbers
$\left\{  \left.  N\left(  \tilde{\psi}_{\#j}\right)  \right\vert j\in\left\{
j_{1H},\cdots,j_{mH}\right\}  \right\}  ,$ produced by the GSO's $\left\{
\left.  \tilde{\psi}_{\#j}\right\vert j\in\left\{  j_{1H},\cdots
,j_{mH}\right\}  \right\}  $, where
\[
\sum_{j\in\left\{  j_{1H},\cdots,j_{mH}\right\}  }N\left(  \tilde{\psi}%
_{\#j}\right)  \kappa\left(  \tilde{\psi}_{\#j}\right)  =\max_{\substack{j\in
\left\{  j_{1},\cdots,j_{m}\right\}  \\N\leq m\leq2s}}\left\{  \sum
_{j_{1},\cdots,j_{m}}N\left(  \tilde{\psi}_{\#j}\right)  \kappa\left(
\tilde{\psi}_{\#j}\right)  \right\}  .
\]
The set of subscripts $\left\{  j_{1},\cdots,j_{m}\right\}  $ are determined
by the constraint
\begin{equation}
\min_{\substack{\left\{  j_{1},\cdots,j_{m}\right\}  \\N\leq m\leq
2s}}\left\vert \sum_{j\in\left\{  j_{1},\cdots,j_{m}\right\}  }N\left(
\tilde{\psi}_{\#j}\right)  -1\right\vert ,\label{con}%
\end{equation}
which cannot be expressed as
\begin{equation}
\sum_{j\in\left\{  j_{1},\cdots,j_{m}\right\}  }N\left(  \tilde{\psi}%
_{\#j}\right)  -1=0
\end{equation}
as there may be no solution. The GSO's $\left\vert \tilde{\psi}_{\#j}%
\right\rangle $, $\left\vert \tilde{\psi}_{\#j+s}\right\rangle $, always have
the same occupation numbers due to the pairing degeneracy, though the energies
$\left\{  \kappa\left(  \tilde{\psi}_{\#j}\right)  ,\kappa\left(  \tilde{\psi
}_{\#j+s}\right)  \right\}  $ might not be the same. The description of the
next most energetic electron can be obtained in an analogues fashion.

\section{Density Matrix Functional\label{DenMatFun}}

In this section we return to the non redundant parameterization of eq.
$\left(  \ref{gem}\right)  ,$ which explicitly involves the phase. As the
FORDO is left invariant by the change of phase, while the $N$-electron state
and energy are not, one can define an AGP DMF (a functional defined over the
domain of AGP states) by%
\begin{equation}
\mathcal{F}\left(  D^{1}\left(  \boldsymbol{c},\boldsymbol{\psi}\right)
\right)  =\min_{\boldsymbol{\theta}}\mathcal{E}\left(  \boldsymbol{\theta
}\right)
\end{equation}
where for fixed $\boldsymbol{c}$ and $\boldsymbol{\psi}$ the energy is just a
function of the phases $\left\{  \left.  \theta_{j}\right\vert 1\leq j<k\leq
s\right\}  $ given by eq. $\left(  \ref{Phases}\right)  $, as
\begin{align}
&  \mathcal{E}\left(  \boldsymbol{\theta}\right)  =\frac{\mathcal{H}\left(
\theta\right)  }{S_{\frac{N}{2}}}=S_{\frac{N}{2}}{}^{-1}\left\{  \sum_{1\leq
j\leq s}n_{j}S_{\frac{N}{2}-1}\left(  \hat{\jmath}\right)  \left\{  \overset
{}{\left\langle \psi_{j}\left\vert h_{1}\psi_{j}\right.  \right\rangle
+\left\langle \psi_{j+s}\left\vert h_{1}\psi_{j+s}\right.  \right\rangle
}+\,\left\langle \psi_{j}\psi_{j+s}\right\vert \left\vert \psi_{j}\psi
_{j+s}\right\rangle \right\}  \right. \nonumber\\
&  +\sum_{1\leq j\neq k\leq s}\left\{  c_{j}c_{k}e^{i\left(  \theta_{j}%
-\theta_{k}\right)  }S_{\frac{N}{2}-1}\left(  \hat{\jmath}\hat{k}\right)
\left\langle \psi_{j}\psi_{j+s}\right\vert \left\vert \psi_{k}\psi
_{k+s}\right\rangle +\frac{1}{2}n_{j}n_{k}S_{\frac{N}{2}-2}\left(  \hat
{\jmath}\hat{k}\right)  \left\langle \psi_{j}\psi_{k}|||\psi_{j}\psi
_{k}\right\rangle \,\overset{}{}\right\}  .
\end{align}


\section{Simple Examples}

In order to make the preceding discussion more tangible one can introduce very
simple examples that involve only two electrons i.e. $N=2$. These demonstrate

\begin{itemize}
\item how one electron orbitals and occupation numbers determine correlated
multi electron (in this case two electron) states and

\item the difference between an OET and a FORDO one, which consequently leads
to the construction of an explicit form for a DMF.
\end{itemize}

These subtle differences can even be observed when one restricts attention to
spin orbitals or even orbitals for the two electron systems. Thus we can
restrict attention to geminals that are low spin eigen functions of $\ $the
spin operator $\mathcal{S}_{0}$ (see Appendix \ref{Spin}) i.e.
\begin{equation}
\mathcal{S}_{0}\left\vert g\right\rangle =0
\end{equation}
formed by spin orbitals $\left\{  \left\vert \varphi_{1}\alpha\right\rangle
,\cdots,\left\vert \varphi_{s}\alpha\right\rangle ,\left\vert \varphi
_{s+1}\beta\right\rangle ,\cdots,\left\vert \varphi_{2s}\beta\right\rangle
\right\}  $. The most general AGP of this type is in fact a geminal of the
form%
\begin{equation}
\left\vert g\left(  \boldsymbol{c},\boldsymbol{\varphi},\boldsymbol{\theta
}\right)  \right\rangle =\sum_{1\leq j\leq s}c_{j}e^{i\theta_{j}}\left\vert
\varphi_{j}\alpha\varphi_{j+s}\beta\right\rangle ,\label{spinorbgem}%
\end{equation}
which describes correlated two electron states if $s\geq2$.

\subsection{Unrestricted Orbital Example}

The simplest example of \emph{different} states that have the \emph{same}
FORDO that are determined by spin orbitals and occupation numbers is given by
eq. $\left(  \ref{spinorbgem}\right)  $ when we set (i)
\begin{equation}
s=2;c_{1}=\frac{1}{\sqrt{2}},c_{2}=\frac{1}{\sqrt{2}};\theta_{1}=0,\theta
_{2}=0
\end{equation}
producing%
\begin{equation}
\left\vert g_{+}\right\rangle =\frac{1}{\sqrt{2}}\left\{  \left\vert
\varphi_{1}\alpha\varphi_{3}\beta\right\rangle +\left\vert \varphi_{2}%
\alpha\varphi_{4}\beta\right\rangle \right\}  \equiv\frac{1}{\sqrt{2}}\left\{
\left\vert 1\alpha3\beta\right\rangle +\left\vert 2\alpha4\beta\right\rangle
\right\}  \label{ex1}%
\end{equation}
and set (ii)
\begin{equation}
s=2;c_{1}=\frac{1}{\sqrt{2}},c_{2}=\frac{1}{\sqrt{2}};\theta_{1}=0,\theta
_{2}=\pi
\end{equation}
producing
\begin{equation}
\left\vert g_{-}\right\rangle =\frac{1}{\sqrt{2}}\left\{  \left\vert
1\alpha3\beta\right\rangle -\left\vert 2\alpha4\beta\right\rangle \right\}  ,
\label{ex2}%
\end{equation}
where the succinct notation
\begin{equation}
j\equiv\varphi_{j};1\leq j\leq2s
\end{equation}
has been introduced. The common FORDO\ is
\begin{equation}
D^{1}=\frac{1}{4}\left\{  \left\vert 1\alpha\right\rangle \left\langle
1\alpha\right\vert +\left\vert 2\alpha\right\rangle \left\langle
2\alpha\right\vert +\left\vert 3\beta\right\rangle \left\langle 3\beta
\right\vert +\left\vert 4\beta\right\rangle \left\langle 4\beta\right\vert
\right\}  \label{D1+-}%
\end{equation}
while the SORDO's are
\begin{equation}
D^{2}\left(  g_{+}\right)  =\frac{1}{2}\left\{  \left\vert 1\alpha
3\beta\right\rangle \left\langle 1\alpha3\beta\right\vert +\left\vert
1\alpha3\beta\right\rangle \left\langle 2\alpha4\beta\right\vert +\left\vert
2\alpha4\beta\right\rangle \left\langle 1\alpha3\beta\right\vert +\left\vert
2\alpha4\beta\right\rangle \left\langle 2\alpha4\beta\right\vert \right\}
\label{D2+}%
\end{equation}
and
\begin{equation}
D^{2}\left(  g_{-}\right)  =\frac{1}{2}\left\{  \left\vert 1\alpha
3\beta\right\rangle \left\langle 1\alpha3\beta\right\vert -\left\vert
1\alpha3\beta\right\rangle \left\langle 2\alpha4\beta\right\vert -\left\vert
2\alpha4\beta\right\rangle \left\langle 1\alpha3\beta\right\vert +\left\vert
2\alpha4\beta\right\rangle \left\langle 2\alpha4\beta\right\vert \right\}  .
\label{D2-}%
\end{equation}
The relationship between the SORDO's eqs. $\left(  \ref{D2+},\ref{D2-}\right)
$ and the general form of the SORDO is shown in Appendix \ref{SORDO}.

In the context of a OET (section \ref{OET}) the states $\left\vert
g_{+}\right\rangle $ and $\left\vert g_{-}\right\rangle $ are determined
by\ the representative spin orbitals and occupation numbers
\begin{equation}
\left\{  \varphi_{1}\alpha,\varphi_{2}\alpha,\varphi_{3}\beta,\varphi_{4}%
\beta\right\}  ;\left\{  \frac{1}{2},\frac{1}{2},\frac{1}{2},\frac{1}%
{2}\right\}  \label{Tilde1}%
\end{equation}
and%
\begin{equation}
\left\{  \varphi_{1}\alpha,-\varphi_{2}\alpha,\varphi_{3}\beta,\varphi
_{4}\beta\right\}  ;\left\{  \frac{1}{2},\frac{1}{2},\frac{1}{2},\frac{1}%
{2}\right\}  \label{Tilde2}%
\end{equation}
respectively. One should note that different representative sets of spin
orbitals could be chosen from the equivalence classes $\left[
\boldsymbol{\psi}\right]  $, but the phase dependence (in these cases the
phases $\theta_{1}=0$ in eq's. $\left(  \ref{Tilde1},\ref{Tilde2}\right)  $,
$\theta_{2}=0$ in eq. $\left(  \ref{Tilde1}\right)  $ and $\theta_{2}=\pi$ eq.
$\left(  \ref{Tilde2}\right)  )$ \emph{have to be included} in the description
of the spin orbital.

These states have different energies which are given by
\begin{align}
&  \mathcal{E}\left(  g_{+}\right)  =\frac{1}{2}\left\{  \overset
{}{\left\langle 1\alpha\left\vert h_{1}1\alpha\right.  \right\rangle
+\left\langle 2\alpha\left\vert h_{1}2\alpha\right.  \right\rangle
+\left\langle 3\beta\left\vert h_{1}3\beta\right.  \right\rangle +\left\langle
4\beta\left\vert h_{1}4\beta\right.  \right\rangle }\right. \nonumber\\
&  \left.  \,+\left\langle 1\alpha3\beta\right\vert \left\vert 1\alpha
3\beta\right\rangle +\left\langle 2\alpha4\beta\right\vert \left\vert
2\alpha4\beta\right\rangle +\left\langle 1\alpha3\beta\right\vert \left\vert
2\alpha4\beta\right\rangle +\left\langle 2\alpha4\beta\right\vert \left\vert
1\alpha3\beta\right\rangle \overset{}{}\right\}  \label{E+}%
\end{align}
and
\begin{align}
&  \mathcal{E}\left(  g_{-}\right)  =\frac{1}{2}\left\{  \overset
{}{\left\langle 1\alpha\left\vert h_{1}1\alpha\right.  \right\rangle
+\left\langle 2\alpha\left\vert h_{1}2\alpha\right.  \right\rangle
+\left\langle 3\beta\left\vert h_{1}3\beta\right.  \right\rangle +\left\langle
4\beta\left\vert h_{1}4\beta\right.  \right\rangle }\right. \nonumber\\
&  \left.  \,+\left\langle 1\alpha3\beta\right\vert \left\vert 1\alpha
3\beta\right\rangle +\left\langle 2\alpha4\beta\right\vert \left\vert
2\alpha4\beta\right\rangle -\left\langle 1\alpha3\beta\right\vert \left\vert
2\alpha4\beta\right\rangle -\left\langle 2\alpha4\beta\right\vert \left\vert
1\alpha3\beta\right\rangle \overset{}{}\right\}  . \label{E-}%
\end{align}


\subsection{Restricted Orbital Example}

The preceding example can be made even simpler by considering restricted spin
orbitals
\begin{align}
\left\vert \varphi_{1}\right\rangle  &  =\left\vert \varphi_{3}\right\rangle
\nonumber\\
\left\vert \varphi_{2}\right\rangle  &  =\left\vert \varphi_{4}\right\rangle
\end{align}
leading to
\begin{align}
\left\vert g_{\text{res}+}\right\rangle  &  =\left(  \sqrt{2}\right)
^{-1}\left\{  \left\vert 1\alpha1\beta\right\rangle +\left\vert 2\alpha
2\beta\right\rangle \right\} \nonumber\\
\left\vert g_{\text{res}-}\right\rangle  &  =\left(  \sqrt{2}\right)
^{-1}\left\{  \left\vert 1\alpha1\beta\right\rangle -\left\vert 2\alpha
2\beta\right\rangle \right\}  . \label{gres}%
\end{align}
The FORDO, SORDO's and energies can be obtained from eqs. $\left(
\ref{D1+-}-\ref{E-}\right)  $ by replacing "3" with "1" and "4" with "2". The
restricted geminals eq. $\left(  \ref{gres}\right)  $ are now pure spin
singlet states and can be factored into a product of a space part and spin
part (see Appendix \ref{Spin}).

Representative spin orbitals and occupation numbers determining the states
$\left\vert g_{\text{res}+}\right\rangle $ and $\left\vert g_{\text{res}%
-}\right\rangle $ are
\begin{align}
&  \left\{  \varphi_{1}\alpha,\varphi_{2}\alpha,\varphi_{1}\beta,\varphi
_{2}\beta\right\}  ;\quad\left\{  \frac{1}{2},\frac{1}{2},\frac{1}{2},\frac
{1}{2}\right\}  ,\\
&  \left\{  \varphi_{1}\alpha,i\varphi_{2}\alpha,\varphi_{1}\beta,i\varphi
_{3}\beta\right\}  ;\quad\left\{  \frac{1}{2},\frac{1}{2},\frac{1}{2},\frac
{1}{2}\right\}
\end{align}
respectively.

\subsection{Density Matrix Functionals\label{DMFEx}}

Consider the manifold of geminals of the form
\begin{equation}
\left\vert g\left(  \boldsymbol{c},\boldsymbol{\psi},\theta\right)
\right\rangle =c_{1}\left\vert 1\alpha3\beta\right\rangle +c_{2}e^{i\theta
}\left\vert 2\alpha4\beta\right\rangle ,
\end{equation}
where
\begin{align}
&  c_{1}^{2}+c_{2}^{2}=1\nonumber\\
&  0\leq\theta<2\pi.
\end{align}
The constraint on the canonical coefficients leads to the parameterization
\begin{align}
c_{1} &  =\cos\chi\nonumber\\
c_{2} &  =\sin\chi\nonumber\\
0 &  \leq\chi\leq\frac{\pi}{2}%
\end{align}
and hence%
\begin{equation}
\left\vert g\left(  \chi,\boldsymbol{\psi},\theta\right)  \right\rangle
=\cos\chi\left\vert 1\alpha3\beta\right\rangle +e^{i\theta}\sin\chi\left\vert
2\alpha4\beta\right\rangle .
\end{equation}
All the states with a fixed $\left\{  \chi,\boldsymbol{\psi}\right\}  $, where
$\boldsymbol{\psi\in}\left[  \boldsymbol{\psi}\right]  $, have a common FORDO
given by%
\[
D^{1}\left(  \chi,\boldsymbol{\psi}\right)  =\cos^{2}\chi\left\{  \left\vert
1\alpha\right\rangle \left\langle 1\alpha\right\vert +\left\vert
3\beta\right\rangle \left\langle 3\beta\right\vert \right\}  +\sin^{2}%
\chi\left\{  \left\vert 2\alpha\right\rangle \left\langle 2\alpha\right\vert
+\left\vert 4\beta\right\rangle \left\langle 4\beta\right\vert \right\}
\]
but different SORDO's%
\begin{align}
D^{2}\left(  \chi,\boldsymbol{\psi}\right)   &  =\cos^{2}\chi\left\vert
1\alpha3\beta\right\rangle \left\langle 1\alpha3\beta\right\vert +\sin^{2}%
\chi\left\vert 2\alpha4\beta\right\rangle \left\langle 2\alpha4\beta
\right\vert +\nonumber\\
&  \frac{1}{2}\left\{  e^{-i\theta}\left\vert 1\alpha3\beta\right\rangle
\left\langle 2\alpha4\beta\right\vert +e^{i\theta}\left\vert 2\alpha
4\beta\right\rangle \left\langle 1\alpha3\beta\right\vert \right\}  \sin2\chi.
\end{align}
These states have different energies that can be expressed by%
\begin{align}
&  \mathcal{E}\left(  \theta\right)  =\left\{  \left\langle 1\alpha\left\vert
h_{1}1\alpha\right.  \right\rangle +\left\langle 3\beta\left\vert h_{1}%
3\beta\right.  \right\rangle \right\}  \cos^{2}\chi+\left\{  \left\langle
2\alpha\left\vert h_{1}2\alpha\right.  \right\rangle +\left\langle
4\beta\left\vert h_{1}4\beta\right.  \right\rangle \right\}  \sin^{2}%
\chi\nonumber\\
&  \qquad\qquad\qquad\qquad\qquad+\left\langle 1\alpha3\beta\right\vert
\left\vert 1\alpha3\beta\right\rangle \cos^{2}\chi+\left\langle 2\alpha
4\beta\right\vert \left\vert 2\alpha4\beta\right\rangle \sin^{2}%
\chi\nonumber\\
&  \qquad\qquad\qquad\qquad\qquad\qquad+\frac{1}{2}\left\{  e^{i\theta
}\left\langle 1\alpha3\beta\right\vert \left\vert 2\alpha4\beta\right\rangle
+e^{-i\theta}\left\langle 2\alpha4\beta\right\vert \left\vert 1\alpha
3\beta\right\rangle \right\}  \sin2\chi.\label{EnTheta}%
\end{align}
After spin integration eq. $\left(  \ref{EnTheta}\right)  $ gives
\begin{align}
\mathcal{E}\left(  \theta\right)   &  =\left\{  \left(  1\left\vert
h_{1}1\right.  \right)  +\left(  3\left\vert h_{1}3\right.  \right)  \right\}
\cos^{2}\chi+\left\{  \left(  2\left\vert h_{1}2\right.  \right)  +\left(
4\left\vert h_{1}4\right.  \right)  \right\}  \sin^{2}\chi\nonumber\\
&  \qquad\qquad\,+\left(  \left.  13\right\vert 13\right)  \cos^{2}%
\chi+\left(  \left.  24\right\vert 24\right)  \sin^{2}\chi+\operatorname{Re}%
\left\{  e^{i\theta}\left(  \left.  13\right\vert 24\right)  \sin
2\chi\right\}  \nonumber\\
&  \equiv k\left(  \chi,\boldsymbol{\psi}\right)  +f\left(  \chi
,\boldsymbol{\psi},\theta\right)  ,
\end{align}
where
\begin{equation}
\left(  \left.  jk\right\vert lm\right)  =\int\varphi_{j}^{\ast}\left(
\boldsymbol{r}_{1}\right)  \varphi_{k}^{\ast}\left(  \boldsymbol{r}%
_{2}\right)  \frac{1}{\left\Vert \boldsymbol{r}_{1}-\boldsymbol{r}%
_{2}\right\Vert }\varphi_{l}\left(  \boldsymbol{r}_{1}\right)  \varphi
_{m}\left(  \boldsymbol{r}_{2}\right)  d\boldsymbol{r}_{1}d\boldsymbol{r}_{2},
\end{equation}
%

\begin{equation}
\left(  j\left\vert h_{1}j\right.  \right)  =\int\varphi_{j}^{\ast}\left(
\boldsymbol{r}\right)  \left\{  -\frac{1}{2}\nabla_{\boldsymbol{r}}%
^{2}+V_{\text{ext}}\left(  \boldsymbol{r}\right)  \right\}  \varphi_{j}\left(
\boldsymbol{r}\right)  d\boldsymbol{r},
\end{equation}%
\begin{align}
k\left(  \chi,\boldsymbol{\psi}\right)   &  =\left\{  \left(  1\left\vert
h_{1}1\right.  \right)  +\left(  3\left\vert h_{1}3\right.  \right)  \right\}
\cos^{2}\chi+\left\{  \left(  2\left\vert h_{1}2\right.  \right)  +\left(
4\left\vert h_{1}4\right.  \right)  \right\}  \sin^{2}\chi\nonumber\\
&  \qquad\qquad\qquad\qquad\qquad\quad+\left(  \left.  13\right\vert
13\right)  \cos^{2}\chi+\left(  \left.  24\right\vert 24\right)  \sin^{2}\chi
\end{align}
and
\begin{equation}
f\left(  \chi,\boldsymbol{\psi},\theta\right)  =\operatorname{Re}\left\{
e^{i\theta}\left(  \left.  13\right\vert 24\right)  \sin2\chi\right\}  .
\end{equation}
Over the domain of such geminals one can define a DMF as
\begin{equation}
\mathcal{F}\left(  \chi,\boldsymbol{\psi}\right)  =\mathcal{F}\left(
D^{1}\left(  \chi,\boldsymbol{\psi}\right)  \right)  =\min_{\boldsymbol{\theta
}}\mathcal{E}\left(  \theta\right)  .
\end{equation}
The only part of the functional $\mathcal{E}$ that depends on $\theta$ is
$f\left(  \chi,\boldsymbol{\psi},\theta\right)  ,$ which one can express as
\begin{equation}
f\left(  \chi,\boldsymbol{\psi},\theta\right)  =\operatorname{Re}\left\{
e^{i\theta}\left(  \left.  13\right\vert 24\right)  \right\}
=\operatorname{Re}\left\{  e^{i\theta}ae^{i\alpha}\right\}  =\operatorname{Re}%
\left\{  e^{i\left(  \theta+\alpha\right)  }a\right\}  =a\cos\left(
\theta+\alpha\right)  ,
\end{equation}
where
\begin{align}
a &  =\left\vert \left(  \left.  13\right\vert 24\right)  \right\vert \\
\alpha &  =\arg\left\{  \left(  \left.  13\right\vert 24\right)  \right\}  .
\end{align}
The minimum of this function can be easily evaluated giving
\begin{align}
f\left(  \chi,\boldsymbol{\psi},\theta_{\min}\right)   &  =\min_{\theta
}\left\{  a\cos\left(  \theta+\alpha\right)  \right\}  =-a\\
\theta_{\min} &  =\pi-\alpha.
\end{align}
Hence the value of the DMF, $\mathcal{F},$ at the point $D^{1}\left(
\chi,\boldsymbol{\psi}\right)  $ is%
\begin{equation}
\mathcal{F}\left(  \chi,\boldsymbol{\psi}\right)  =k\left(  \chi
,\boldsymbol{\psi}\right)  -2\left\vert \left(  \left.  13\right\vert
24\right)  \right\vert \sin2\chi
\end{equation}
and the minimum of the AGP energy functional $\mathcal{E}\left(
\boldsymbol{c}\left(  \chi\right)  ,\boldsymbol{\psi},\theta\right)  ,$ which
determines the AGP\ ground state, is given by
\begin{equation}
\mathcal{E}\left(  \left\vert g_{GS}\right\rangle \right)  =\min
_{\chi,\boldsymbol{\psi}}\mathcal{F}\left(  \chi,\boldsymbol{\psi}\right)  .
\end{equation}


\section{Discussion}

We have shown how the universally used orbital picture of electrons, that is
only valid for uncorrelated systems, can be extended to a description of
correlated ones. This allows one to maintain the single electron description
of a mutli electron system and construct a pictorial model as outline in
subsection \ref{OrbModel}. The manifold of $N$-electron states that are
determined by GSO's and their occupation numbers are formed by AGP\ states,
which describe systems that contain an even number of electrons. In order to
describe systems with an odd number of electrons by a OET one can use
Generalized AGP (GAGP) states \cite{1966A03263} that are formed by taking an
antisymmetrized product of a GSO\ with an AGP state. AGP\ states are not in
general ground states for molecular systems, but can produce approximate
ground states that have a lower energy than IPS's.

OET's are more general than FORDO ones as GSO's with the same occupation
numbers can produce different $N$-electron states while producing the same
FORDO. We constructed some simple two electron examples that displayed this
feature, which we subsequently used to construct an example of an explicit DMF
over the domain of geminals.

\section{Acknowledgements}

I\ would like to thank Dr. Vince Ortiz for valuable discussions on this topic.

\appendix


\section{Representative GSO Pairs\label{GSOpairs}}

One can expand a geminal in a form that explicitly displays degeneracies as
\begin{equation}
\left\vert g\left(  \boldsymbol{c},\boldsymbol{\psi},\boldsymbol{\theta
}\right)  \right\rangle =\sum_{1\leq k\leq d}c_{j_{k}}\sum_{s_{k}+1\leq
j_{k}\leq s_{k}+d_{k}}e^{i\theta_{j_{k}}}\left\vert \psi_{j_{k}}\psi_{j_{k}%
+s}\right\rangle ,
\end{equation}
where $\left\{  d_{1},\cdots,d_{d}\right\}  $ are the degree of degeneracies
of the canonical coefficients $\left\{  c_{1},\cdots,c_{d}\right\}  ,$
$s_{k}=\sum\limits_{1\leq l\leq k}d_{l}$ and $\sum\limits_{1\leq k\leq d}%
d_{k}=s_{d}=s.$ The simplest example of a degenerate geminal can be realized
by setting $s=3,d=2,\left\{  d_{1},d_{2}\right\}  =\left\{  2,1\right\}  ,$
i.e.
\begin{equation}
\left\vert g\left(  \boldsymbol{c},\boldsymbol{\psi},\boldsymbol{\theta
}\right)  \right\rangle =c_{1}\left\{  e^{i\theta_{1}}\left\vert \psi_{1}%
\psi_{4}\right\rangle +e^{i\theta_{2}}\left\vert \psi_{2}\psi_{5}\right\rangle
\right\}  +c_{2}e^{i\theta_{3}}\left\vert \psi_{3}\psi_{6}\right\rangle .
\end{equation}
Following \cite{OrtizWeiner2006,Weiner83a,AGPCoherent} one can define a set of
operators associated with the degenerate canonical coefficient $c_{1}$ as%
\begin{equation}
\left(
\begin{array}
[c]{c}%
u_{1}\\
u_{2}\\
u_{3}\\
u_{4}\\
v_{1}\\
v_{2}\\
v_{3}\\
v_{4}%
\end{array}
\right)  =\left(
\begin{array}
[c]{cccccccc}%
e^{-i\theta_{1}} & 0 & 0 & 0 & 0 & 0 & 0 & e^{-i\theta_{2}}\\
0 & e^{-i\theta_{1}} & 0 & 0 & 0 & 0 & -e^{-i\theta_{2}} & 0\\
0 & 0 & e^{-i\theta_{1}} & 0 & 0 & -e^{-i\theta_{2}} & 0 & 0\\
0 & 0 & 0 & e^{-i\theta_{1}} & e^{-i\theta_{2}} & 0 & 0 & 0\\
0 & 0 & 0 & -e^{i\theta_{2}} & e^{i\theta_{1}} & 0 & 0 & 0\\
0 & 0 & e^{i\theta_{2}} & 0 & 0 & e^{i\theta_{1}} & 0 & 0\\
0 & e^{i\theta_{2}} & 0 & 0 & 0 & 0 & e^{i\theta_{1}} & 0\\
-e^{i\theta_{2}} & 0 & 0 & 0 & 0 & 0 & 0 & e^{i\theta_{1}}%
\end{array}
\right)  \left(
\begin{array}
[c]{c}%
a_{1}^{\dagger}a_{2}\\
a_{1}^{\dagger}a_{5}\\
a_{4}^{\dagger}a_{2}\\
a_{4}^{\dagger}a_{5}\\
a_{2}^{\dagger}a_{1}\\
a_{5}^{\dagger}a_{1}\\
a_{2}^{\dagger}a_{4}\\
a_{5}^{\dagger}a_{4}%
\end{array}
\right)  ,
\end{equation}
where the operators $\boldsymbol{u}=\left\{  u_{1},u_{2},u_{3},u_{4}\right\}
,\boldsymbol{v}=\left\{  v_{1},v_{2},v_{3},v_{4}\right\}  $ have the
properties
\begin{equation}
v\left\vert g\right\rangle =0,\quad u\left\vert g\right\rangle \neq0;\quad
u\in\boldsymbol{u},v\in\boldsymbol{v.}%
\end{equation}
Using the set $\boldsymbol{v}$ one can form a basis, $\boldsymbol{v}^{\prime}%
$, of anti self adjoint operators
\begin{align}
v_{1}^{\prime}  & =i\left(  v_{1}+v_{1}^{\dagger}\right)  \nonumber\\
v_{2}^{\prime}  & =v_{1}-v_{1}^{\dagger}\nonumber\\
v_{3}^{\prime}  & =i\left(  v_{2}+v_{2}^{\dagger}\right)  \nonumber\\
v_{4}^{\prime}  & =v_{2}-v_{2}^{\dagger}%
\end{align}
that together with the operators
\begin{align}
&  a_{1}^{\dagger}a_{4}-a_{4}^{\dagger}a_{1},i\left(  a_{1}^{\dagger}%
a_{4}+a_{4}^{\dagger}a_{1}\right)  ,i\left(  a_{1}^{\dagger}a_{1}%
-a_{4}^{\dagger}a_{4}\right)  ,\nonumber\\
&  a_{2}^{\dagger}a_{5}-a_{5}^{\dagger}a_{2},i\left(  a_{2}^{\dagger}%
a_{5}+a_{5}^{\dagger}a_{2}\right)  ,i\left(  a_{2}^{\dagger}a_{2}%
-a_{5}^{\dagger}a_{5}\right)
\end{align}
form a representation of the unitary symplectic lie algebra $usp\left(
4\right)  $ \cite{Gilmore_book}, which produce a representation of the unitary
symplectic group $USp\left(  4\right)  $ \cite{Gilmore_book} over the subspace
generated by $\left\{  \psi_{1},\psi_{2},\psi_{4},\psi_{5}\right\}  $, that
leaves $\left\vert g\right\rangle $ invariant. Thus any GSO produced by the
action of this representation of $USp\left(  4\right)  $ produces the same
geminal, hence all the GSO's $\left\{  USp\left(  4\right)  \left\vert
\boldsymbol{\psi}\right\rangle \right\}  $ belong to the same equivalence
class $\left[  \boldsymbol{\psi}\right]  ,$ which generalizes the equivalence
relationship given in eq. $\left(  \ref{Equival}\right)  .$ One can always
form sets $\boldsymbol{v}^{\prime}$ even if no degeneracy exists, however in
these cases they do not belong to representations of unitary Lie algebras only
symplectic ones. The set $\boldsymbol{u}^{\prime}$, formed in an analogues
manner to the set $\boldsymbol{v}^{\prime}$, is contained in a representation
of the lie algebra $u\left(  4\right)  $ that produces a representation of the
group $U\left(  4\right)  $ that changes $\left\vert g\right\rangle .$

\section{Spin\label{Spin}}

The spin squared operator, $\mathcal{S}^{2}$, is given by
\begin{equation}
\mathcal{S}^{2}=\mathcal{S}_{-}\mathcal{S}_{+}+\mathcal{S}_{0}^{2}%
+\mathcal{S}_{0},
\end{equation}
where
\begin{align}
\mathcal{S}_{+} &  =\sum_{1\leq j\leq s}a_{j\alpha}^{\dagger}a_{j\beta},\\
\mathcal{S}_{-} &  =\sum_{1\leq j\leq s}a_{j\beta}^{\dagger}a_{j\alpha}%
=S_{+}^{\dagger},\\
\mathcal{S}_{0} &  =\sum_{1\leq j\leq s}a_{j\alpha}^{\dagger}a_{j\alpha}%
-\sum_{1\leq j\leq s}a_{j\beta}^{\dagger}a_{j\beta}.
\end{align}
and the fermion second quantized operators $\left\{  \left.  a_{j\sigma
}^{\dagger}\right\vert 1\leq j\leq s;\sigma=\alpha,\beta\right\}  $ are
defined by%
\begin{align}
a_{j\sigma}^{\dagger}\left\vert \phi\right\rangle  &  =\left\vert \varphi
_{j}\sigma\right\rangle \nonumber\\
\left[  a_{j\sigma}^{\dagger},a_{k\nu}\right]  _{+} &  =a_{j\sigma}^{\dagger
}a_{k\nu}+a_{k\nu}a_{j\sigma}^{\dagger}=\delta_{jk}\delta_{\sigma\nu}.
\end{align}
The action of $\mathcal{S}^{2}$ on $\left\vert g_{\pm}\right\rangle $ is given
by
\begin{align}
\mathcal{S}^{2}\left\vert g_{\pm}\right\rangle  &  =\mathcal{S}^{2}\left(
\sqrt{2}\right)  ^{-1}\left\{  \left\vert \varphi_{1}\alpha\varphi_{3}%
\beta\right\rangle \pm\left\vert \varphi_{2}\alpha\varphi_{4}\beta
\right\rangle \right\}  =\mathcal{S}_{-}\mathcal{S}_{+}\left(  \sqrt
{2}\right)  ^{-1}\left\{  \left\vert \varphi_{1}\alpha\varphi_{3}%
\beta\right\rangle \pm\left\vert \varphi_{2}\alpha\varphi_{4}\beta
\right\rangle \right\}  \nonumber\\
&  =\mathcal{S}_{-}\left(  \sqrt{2}\right)  ^{-1}\left\{  \left\vert
\varphi_{1}\alpha\varphi_{3}\alpha\right\rangle \pm\left\vert \varphi
_{2}\alpha\varphi_{4}\alpha\right\rangle \right\}  \nonumber\\
&  =\left(  \sqrt{2}\right)  ^{-1}\left\{  \left\vert \varphi_{1}\beta
\varphi_{3}\alpha\right\rangle +\left\vert \varphi_{1}\alpha\varphi_{3}%
\beta\right\rangle \pm\left\vert \varphi_{2}\beta\varphi_{4}\alpha
\right\rangle \pm\left\vert \varphi_{2}\alpha\varphi_{4}\beta\right\rangle
\right\}  ,
\end{align}
which displays that $\left\vert g_{\pm}\right\rangle $ are not in general
eigen states of $\mathcal{S}^{2}$. However if one sets $\varphi_{3}%
=\varphi_{1}$ and $\varphi_{4}=\varphi_{2}$ to form restricted geminals one
obtains
\begin{equation}
\mathcal{S}^{2}\left\vert g_{\text{res},\pm}\right\rangle =0,
\end{equation}
showing that $\left\vert g_{\pm}\right\rangle $ are low spin singlet eigen
states of $\mathcal{S}^{2}$ that factor into a product of space and spin parts
i.e.
\begin{equation}
\left\vert g_{\text{res},\pm}\right\rangle =\frac{1}{\sqrt{2}}\left\{
\left\{  \varphi_{1}\left(  1\right)  \otimes\varphi_{1}\left(  2\right)
\pm\varphi_{2}\left(  1\right)  \otimes\varphi_{2}\left(  2\right)  \right\}
\left\{  \alpha\left(  1\right)  \otimes\beta\left(  2\right)  -\beta\left(
1\right)  \otimes\alpha\left(  2\right)  \right\}  \right\}  .
\end{equation}


GSO's, $\left\vert \psi\right\rangle $, are not eigen states of $\mathcal{S}%
_{0}$ as they are linear combinations of both alpha and beta spin components
$\left\{  \left.  \left\vert \varphi_{j}\alpha\right\rangle ,\left\vert
\varphi_{j}\beta\right\rangle \right\vert 1\leq j\leq s\right\}  $ i.e.%
\begin{equation}
\left\vert \psi\right\rangle =\sum_{1\leq j\leq s}\left\{  a_{j}\left\vert
\varphi_{j}\alpha\right\rangle +b_{j}\left\vert \varphi_{j}\beta\right\rangle
\right\}  ,
\end{equation}
which individually are eigen states of $\mathcal{S}_{0}$ that satisfy
\begin{equation}
\mathcal{S}_{0}\left\vert \varphi_{j}\alpha\right\rangle =\frac{1}%
{2}\left\vert \varphi_{j}\alpha\right\rangle ,\quad\mathcal{S}_{0}\left\vert
\varphi_{j}\beta\right\rangle =-\frac{1}{2}\left\vert \varphi_{j}%
\beta\right\rangle .
\end{equation}


\section{SORDO\label{SORDO}}

Consider the SORDO of the geminal $\left\vert g_{-}\right\rangle $
\begin{equation}
D^{2}\left(  g_{-}\right)  =\frac{1}{2}\left\{  \left\vert 1\alpha
3\beta\right\rangle \left\langle 1\alpha3\beta\right\vert -\left\vert
1\alpha3\beta\right\rangle \left\langle 2\alpha4\beta\right\vert -\left\vert
2\alpha4\beta\right\rangle \left\langle 1\alpha3\beta\right\vert +\left\vert
2\alpha4\beta\right\rangle \left\langle 2\alpha4\beta\right\vert \right\}
,\label{Dirgm}%
\end{equation}
where
\begin{align}
m &  =\frac{N}{2}=\frac{2}{2}=1,\\
s &  =2,\\
n_{1} &  =n_{2}=\frac{1}{2},\\
c_{1} &  =c_{2}=\frac{1}{\sqrt{2}},\\
\theta_{1} &  =0,\quad\theta_{2}=\pi.
\end{align}
Substituting for $m,s$ and $r$ in the general form for the SORDO results in
\begin{align}
&  D^{2}\left(  g_{-}\right)  =\left(  S_{1}\right)  ^{-1}\left\{  \sum_{1\leq
j\leq2}n_{j}S_{0}\left(  \hat{\jmath}\right)  \left\vert \psi_{j}\psi
_{j+s}\right\rangle \left\langle \psi_{j}\psi_{j+s}\right\vert \right.
\nonumber\\
&  +\sum_{1\leq j\neq k\leq2}c_{j}c_{k}e^{i\left(  \theta_{j}-\theta
_{k}\right)  }S_{0}\left(  \hat{\jmath}\hat{k}\right)  \left\vert \psi_{j}%
\psi_{j+s}\right\rangle \left\langle \psi_{k}\psi_{k+s}\right\vert \left.
+\sum_{\substack{1\leq j<k\leq4\\k\neq j+2}}n_{j}n_{k}S_{-1}\left(
\hat{\jmath}\hat{k}\right)  \left\vert \psi_{j}\psi_{k}\right\rangle
\left\langle \psi_{j}\psi_{k}\right\vert \right\}  .
\end{align}
Noting that the symmetric functions are given by
\begin{align}
S_{1} &  =n_{1}+n_{2}=1\\
S_{0}\left(  \hat{1}\hat{2}\right)   &  =1\\
S_{0}\left(  \hat{1}\right)   &  =S_{0}\left(  \hat{2}\right)  =1\\
S_{-1}\left(  \hat{1}\hat{2}\right)   &  =0.
\end{align}
and substituting for $\boldsymbol{c},\boldsymbol{n}$ and $\boldsymbol{\theta}$
one obtains
\begin{align}
&  D^{2}\left(  g_{-}\right)  =\frac{1}{2}\left\{  \left\vert \psi_{1}\psi
_{3}\right\rangle \left\langle \psi_{1}\psi_{3}\right\vert +\left\vert
\psi_{2}\psi_{4}\right\rangle \left\langle \psi_{2}\psi_{4}\right\vert
\right\}  \nonumber\\
&  \qquad\qquad\qquad\qquad\qquad\qquad+\frac{1}{2}\left\{  \left\vert
\psi_{1}\psi_{3}\right\rangle \left\langle \psi_{2}\psi_{4}\right\vert
e^{-i\pi}+\left\vert \psi_{2}\psi_{4}\right\rangle \left\langle \psi_{1}%
\psi_{3}\right\vert e^{i\pi}\right\}  \\
&  =\frac{1}{2}\left\{  \left\vert \psi_{1}\psi_{3}\right\rangle \left\langle
\psi_{1}\psi_{3}\right\vert -\left\vert \psi_{1}\psi_{3}\right\rangle
\left\langle \psi_{2}\psi_{4}\right\vert -\left\vert \psi_{2}\psi
_{4}\right\rangle \left\langle \psi_{1}\psi_{3}\right\vert +\left\vert
\psi_{2}\psi_{4}\right\rangle \left\langle \psi_{2}\psi_{4}\right\vert
\right\}
\end{align}
and finally substituting for $\boldsymbol{\psi}$ one gets eq. $\left(
\ref{Dirgm}\right)  .$

\bibliographystyle{apsrev}
\bibliography{agpopt,bw,novel,San2008,coleman,Coleman2,bibrefs}

\end{document}